\section{模型检验、敏感性与适用边界}

\subsection{数值一致性和独立回放}

计算链分为“优化/决策器”和“规则模拟器”两个入口。问题一由MILP产生行动变量，再由不参与优化的逐日模拟器检查；第三关Bellman值与1024条全序列加权回放一致；第五关均衡计划由独立的路径评价器重算资源和收益；第四、六关所有Monte Carlo样本均逐日执行同一规则内核。当前20项自动测试覆盖地图连通与对称、沙暴禁行、非矿山挖矿拒绝、起点超重拒绝、村庄行动前购买、已知天气参考流水、在线天气复现和多人同边消耗。

\subsection{训练—测试隔离与幸存者偏差}

第四关训练种子为20260721，最终测试种子从20260821开始；第六关装载训练使用与最终20260921测试不同的随机流。候选选择只读取训练结果，测试输出不反向修改风险预算。设成功样本终值为 $J_j$，失败指示为 $F_j$，本文的总体估计量为
\begin{equation}
\widehat{\mu}=\frac1N\sum_{j=1}^N (1-F_j)J_j,
\end{equation}
而不是只在 $F_j=0$ 的子样本上取均值。故高收益、高失败的策略不会因剔除失败局而显得虚假优越。失败率区间采用Wilson得分区间\cite{wilson1927}。

\subsection{消融：为什么不能复制单人策略}

第六关镜像挖矿是对“忽略交互项”的直接消融。单人第四关B5在基准天气总体均值10162.9元、失败率0.055\%；三人原样复制后，基准失败率升至97.865\%，总体均值降到142.6元。变化并非天气或地图造成，而是同边移动从2倍升为6倍。进一步从镜像改为分路但保留矿山竞争，失败率降至12.362\%；再彻底放弃矿山并分路直达，失败率降至0.0167\%。这一逐层消融把外部性来源从相关性转化为可解释机制。

\subsection{六关证据矩阵}

\begin{figure}[H]
  \centering
  \includegraphics[width=0.88\textwidth]{evidence_matrix.pdf}
  \caption{六关核心结果与结论证据等级}
  \label{fig:evidence}
\end{figure}

\subsection{模型优点}

\begin{enumerate}
  \item \textbf{信息结构一致。} 三类模型都明确天气在行动前或行动后何时可见，杜绝未来信息泄漏；随机规划中的逐期信息适应原则与多阶段随机优化框架一致\cite{shapiro2014}。
  \item \textbf{结论强度可审计。} 全局最优、全枚举、最佳响应和样本外推荐使用不同证据标签，评审可快速识别哪些结论依赖概率假设。
  \item \textbf{失败被纳入效用。} 选模同时控制均值和失败风险，避免只比较到达者。
  \item \textbf{一般性强。} 地图、天气、资源与多人系数通过场景对象注入；更换附件无需重写规则内核。
  \item \textbf{机制可解释。} 同边玩家日、同矿玩家日、库存和现金流水把结果落回题面规则，而非黑箱给数。
\end{enumerate}

\subsection{局限与可证伪条件}

\begin{boundarybox}
\textbf{分布边界：}题面只说“较少沙暴”，没有给出概率或相关性。若真实沙暴概率、连续沙暴聚集程度或高温—沙暴转移显著高于本文不利分布，第四、六关风险率需重新估计。

\textbf{均衡边界：}第五关完成了有限策略空间的完整最佳响应，结论可称Nash均衡；第六关只比较一组可解释反馈策略，没有证明Markov完美均衡。

\textbf{行为边界：}公开随机角色分配实现事前公平，但题面没有强制协调机制。若玩家拒绝角色约定，应使用非合作动态博弈重新求解。

\textbf{数据边界：}地图由附件图形人工重建并经拓扑测试；若原图邻接判定有误，最短路与全部结果都会改变。节点邻接表因此作为可审计数据随代码交付。
\end{boundarybox}

最可能推翻推荐的三类新信息是：官方给出天气转移矩阵；第六关允许有约束力的合作和转移支付；附件地图的邻接边被官方修正。任一变化发生时，应保留规则模拟器，仅替换天气模型、博弈机制或邻接表后重跑。
